\selectlanguage{brazil}%
Modelo 
\begin{align}
\dot{\positionVector}^{\navigationSystem} & =-_{\bodySystem}^{\navigationSystem}\matrixRotation\velocityVector^{\bodySystem},\label{eq:-11}\\
\dot{\eulerAngleVector} & =-\matrixTransformationOfKinematics\angularRateVector,\label{eq:-12}\\
\mass\dot{\velocityVector}^{\bodySystem} & =\angularRateVector\times\velocityVector^{\bodySystem}-{}_{\navigationSystem}^{\bodySystem}\matrixRotation\gravityVector^{\navigationSystem}-\forceVectorAerodynamic[^{\bodySystem}],\label{eq:-9}\\
\tensorOfInertiaMatrix\dot{\angularRateVector}^{\bodySystem} & =\angularRateVector^{\bodySystem}\timesOperator\tensorOfInertiaMatrix\angularRateVector^{\bodySystem}-\momentVector^{\bodySystem}.\label{eq:-10-1}
\end{align}

The state, derivative, and control vectors for the system are defined
as: 
\begin{align}
\stateVector & =\left[\begin{array}{cccc}
\positionVector^{\navigationSystem} & \eulerAngleVector & \velocityVector^{\bodySystem} & \angularRateVector\end{array}\right],\label{eq:l-17}\\
\dot{\stateVector} & =\left[\begin{array}{cccc}
\dot{\positionVector}^{\navigationSystem} & \dot{\eulerAngleVector} & \dot{\velocityVector}^{\bodySystem} & \dot{\angularRateVector}\end{array}\right],\label{eq:l-18}\\
\controlVector & =\left[\begin{array}{cccc}
\forcePropulsion & \rightElevon & \leftElevon & \rudder\end{array}\right]\label{eq:l-45}
\end{align}
These equations contain terms for aerodynamic forces ($\forceVectorAerodynamic$)
and moments ($\momentVector^{\bodySystem}$), which are functions
of the vehicle's state and control inputs.

No caso do planador, não exite $\forcePropulsion$ nem $\rudder$.
Controle de um planador.

The force and moment coefficients are then expressed as linear functions
of the perturbations in these aerodynamic states and control surface
deflections: \begin{subequations} 
\begin{align}
\coefficientBodyForceAxial & =\coefficientBodyForceAxialZeroLift+\coefficientBodyForceAxialIncrementForAngleAtack\angleAttack+\coefficientBodyForceAxialIncrementForRightElevon\rightElevon+\coefficientBodyForceAxialIncrementForLeftElevon\leftElevon+\coefficientBodyForceAxialIncrementForAngularRateq\frac{\angularRateq\referenceLongitudinalLength}{2\velocity}.\label{eq:l}\\
\coefficientLateralForceWind & =\coefficientLateralForceWindIncrementForAngleSideSlip\angleSideslip+\coefficientLateralForceWindIncrementForDotAngleSideSlip\frac{\dot{\angleSideslip}\referenceLateralLength}{2\velocity}+\coefficientLateralForceWindIncrementForAngularRatep\frac{\angularRatep\referenceLateralLength}{2\velocity}+\coefficientLateralForceWindIncrementForAngularRater\frac{\angularRater\referenceLateralLength}{2\velocity},\label{eq:l-3}\\
\coefficientBodyForceNormal & =\coefficientBodyForceNormalZeroLift+\coefficientBodyForceNormalIncrementForAngleAtack\angleAttack+\coefficientBodyForceNormalIncrementForRightElevon\rightElevon+\coefficientBodyForceNormalIncrementForLeftElevon\leftElevon+\coefficientBodyForceNormalIncrementForAngularRateq\frac{\angularRateq\referenceLongitudinalLength}{2\velocity}+\coefficientBodyForceNormalIncrementForDotAngleAtack\frac{\dot{\angleAttack}\referenceLongitudinalLength}{2\velocity},\label{eq:l-1}
\end{align}
\end{subequations}\begin{subequations} 
\begin{align}
\coefficientPitchingMoment & =\coefficientPitchingMomentZeroLift+\coefficientPitchingMomentIncrementForAngleAtack\angleAttack+\coefficientPitchingMomentIncrementForRightElevon\rightElevon+\coefficientPitchingMomentIncrementForLeftElevon\leftElevon+\coefficientPitchingMomentIncrementForAngularRateq\frac{\angularRateq\,\referenceLongitudinalLength}{2\velocity}+\coefficientPitchingMomentIncrementForDotAngleAtack\frac{\dot{{\angleAttack}}\referenceLongitudinalLength}{2\velocity}.\label{eq:l-2}\\
\coefficientRollingMoment & =\coefficientRollingMomentIncrementForAngleSideSlip\angleSideslip+\coefficientRollingMomentIncrementForLeftElevon\leftElevon+\coefficientPitchingMomentIncrementForRightElevon\rightElevon+\coefficientRollingMomentIncrementForAngularRatep\frac{\angularRatep\,\referenceLateralLength}{2\velocity}+\coefficientRollingMomentIncrementForAngularRater\frac{\angularRater\,\referenceLateralLength}{2\velocity}+\coefficientRollingMomentIncrementForDotAngleSideSlip\frac{\dot{{\angleSideslip}}\referenceLongitudinalLength}{2\velocity},\label{eq:-2}\\
\coefficientYawingMoment & =\coefficientYawingMomentIncrementForAngleSideSlip\angleSideslip+\coefficientYawingMomentIncrementForLeftElevon\leftElevon+\coefficientPitchingMomentIncrementForRightElevon\rightElevon+\coefficientYawingMomentIncrementForAngularRatep\frac{\angularRatep\,\referenceLateralLength}{2\velocity}+\coefficientYawingMomentIncrementForAngularRater\frac{\angularRater\,\referenceLateralLength}{2\velocity}+\coefficientYawingMomentIncrementForDotAngleSideSlip\frac{\dot{{\angleSideslip}}\referenceLateralLength}{2\velocity}.\label{eq:-6}
\end{align}

\par \end{subequations}

The aerodynamic forces and moments are modeled using a coefficient-based
approach. The total force and moment vectors are defined as: 
\begin{align}
\forceVectorAerodynamic & =\dynamicPressure\referenceArea\begin{bmatrix}-\coefficientBodyForceAxial\\
\coefficientLateralForceWind\\
-\coefficientBodyForceNormal
\end{bmatrix}+\left[\begin{array}{c}
\forcePropulsion\\
0\\
0
\end{array}\right].\label{eq:-7}
\end{align}
\begin{align}
\momentVector^{\bodySystem} & =\dynamicPressure\referenceArea\begin{bmatrix}\referenceLateralLength\coefficientRollingMoment\\
\referenceLongitudinalLength\coefficientPitchingMoment\\
\referenceLateralLength\coefficientYawingMoment
\end{bmatrix}\label{eq:-8}
\end{align}

As entradas de controle são dois elevons, um direito e outro esquerdo
$\leftElevon$ and $\rightElevon$. Para isso, será implementado duas
malhas. A condição de equilíbrio é obtida através de um algorítimo
de TRIM. No equilíbrio de um planador as condições são: 
\begin{align}
\angleAttack & =\eulerTheta-\angleFlightPath,\\
\angleSideslip & =0\rightarrow\velocityv=0,\\
\angularRateVector & =0,\\
\eulerPhi & =0,\\
\eulerPsi & =\text{não importa agora, pode ser zero},
\end{align}

\begin{enumerate}
\item Malha translacional 
\begin{enumerate}
\item defino os estados de posição ($\positionVector$) e velocidade ($\velocityVector$)
de translação no sistema \gls{ecef}. A saída do simulador compõe
os estados atuais, ou seja no instante $k$. 
\item Os alvos são $\positionVector_{s}$ e $\velocityVector_{s}$. O alvo
de velocidade é obtido a partir do número de Mach $\machNumber$,
condição de não derrapagem (como $\angleSideslip$ deve ser nulo,
a velocidade no corpo lateral deve ser nula $\velocityv=0$). Com
$\machNumber$ e ângulo de ataque ($\angleAttack$ obtido por um algorítimo
de TRIM) as velocidades no corpo são calculadas. Primeira dúvida,
e a referencia de posição? eu integro a velocidade desejada? 
\item A lei de controle considera os erros $\mathbf{e}_{\positionVector}=\positionVector-\positionVector_{s}$,
$\mathbf{e_{\velocityVector}}=\velocityVector-\velocityVector_{s}$
e $\mathbf{e}_{\accelerationVector}=\accelerationVector-\accelerationVector_{s}$,
ou, por alocação de polos
\[
\]
\end{enumerate}
\item Malha rotacional 
\begin{enumerate}
\item d \selectlanguage{american}%
\end{enumerate}
\end{enumerate}

